\section*{Abstract}
%Much of extant genetic material is known to have been produced by genetic duplication events.
Gene duplication is widely recognized as a source of evolutionary novelty and a key factor in the evolution of organismal complexity.
Gene duplicates provide primed genetic material, which can produce novel adaptations (i.e., neofunctionalization) or partition ancestral functions among duplicates (i.e., subfunctionalization).
However, %across the tree of life, no ready counterfactual exists to directly test the aggregate impact of gene duplication on complexity over long evolutionary timescales.
its contribution to the evolution of complexity is difficult to disentangle from the effects of genome expansion alone.
Here, we use the Avida Digital Evolution platform to investigate whether duplication boosts evolved structural complexity (measured as the number of adaptive genome sites) or functional complexity (the sophistication of evolved traits)?
If so, does evolved complexity rely directly on duplicated genetic information?
% In each case, we compare with controls that manipulate the qualities of the inserted sequences.
Duplications lead to the evolution of virtual organisms with higher complexity, in both measures, even when compared to controls with equivalent-sized insertions of random or neutral generic material.
However, the causal role of genetic information differed between measures of complexity.
Duplications increased structural complexity by growing genome size, on average, larger than randomized and empty-insert controls.
By contrast, duplicated genetic information more weakly affected coding-site structure, with complexity here attributable to larger genome size alone.
Increases in genome length sufficed to drive evolution of \textit{simple} phenotypic traits, but duplicated genetic information, in specific, contributed to the evolution of \textit{complex} phenotypic traits.
% Duplications also increased functional complexity by accelerating evolutionary innovation compared to controls.
% Such complex traits arose disproportionately from duplicated coding sites, and seldom arose in large-genome-only treatments.
In sum, we show that while the duplication of genetic information can drive genome expansion and complex adaptation, simpler adaptations and structural complexity may follow from larger genome size alone.
Notably, this finding integrates adaptationist and neutralist framings of complexity, as duplicated information actively catalyzes complex functionality while some aspects of structural complexity may accrue passively.
